% !TEX program = pdflatex
\documentclass[11pt]{article}
\usepackage[a4paper,margin=1in]{geometry}
\usepackage{amsmath,amssymb,amsfonts}
%\usepackage{enumitem}
\usepackage{multicol}
\usepackage{booktabs}
%\usepackage{titlesec}
\usepackage{hyperref}
\hypersetup{colorlinks=true,linkcolor=black,urlcolor=blue}
%\setlist[enumerate]{itemsep=.4em,topsep=.4em}
\setlength{\parskip}{.5em}
\setlength{\parindent}{0pt}

\title{GED Algebra Practice: Algebra Basics, Expressions, and Polynomials}
\author{MacLachlan College}
\date{\today}

\begin{document}
\maketitle
\begin{center}
This worksheet contains 20 practice problems aligned to GED Algebra fundamentals. \textbf{Do not} use a calculator unless instructed by your teacher. Show all work.
\end{center}

\hrulefill

\section*{Part A: Problems}
\textit{Work each problem. Full solutions appear on page~\pageref{sec:solutions}.}

\begin{enumerate}
% Algebra Basics (1–7)
\item Solve for $x$: $x - 9 = 14$.
\item Solve for $x$: $5x + 7 = 2$.
\item Solve for $x$: $\dfrac{x}{6} - 3 = 1$.
\item Solve for $x$: $-4(2 - x) = 12$.
\item A number decreased by 15 is equal to twice the number minus 27. Find the number.
\item The perimeter of a rectangle is $54$ cm. The length is $3$ cm more than twice the width. Find the width.
\item Solve the system by reasoning: If $3$ more than a number is $20$, what is twice the number?

% Expressions (8–13)
\item Simplify completely: $7x - 3 + 5x + 8$.
\item Simplify: $2(3x - 5) - (x + 7)$.
\item Simplify: $6 - [2 - (3x - 4)]$.
\item Evaluate $4y - 2(3y - 5)$ when $y = -2$.
\item Given $a=3$ and $b=-5$, compute $2a^2 - ab + b$.
\item Write an equivalent expression by factoring: $12x + 18$.

% Polynomials (14–20)
\item Add: $(3x^2 - 2x + 5) + (x^2 + 7x - 3)$.
\item Subtract: $(5x^2 + 6x - 4) - (2x^2 - x + 9)$.
\item Multiply: $-3x(4x - 7)$.
\item Expand: $(x + 6)(x - 2)$.
\item Expand: $(2x - 3)(x + 5)$.
\item Factor completely: $x^2 + 9x + 14$.
\item Factor completely: $6x^2 + 11x + 3$.
\end{enumerate}

\newpage
\section*{Part B: Full Solutions}
\label{sec:solutions}
\begin{enumerate}
\item $x - 9 = 14$ \\ Add $9$ to both sides: $x = 23$.

\item $5x + 7 = 2$ \\ Subtract $7$: $5x = -5$. Divide by $5$: $x = -1$.

\item $\dfrac{x}{6} - 3 = 1$ \\ Add $3$: $\dfrac{x}{6} = 4$. Multiply by $6$: $x = 24$.

\item $-4(2 - x) = 12$ \\ Distribute: $-8 + 4x = 12$. Add $8$: $4x = 20$. Divide by $4$: $x = 5$.

\item Let the number be $n$. "A number decreased by 15" is $n - 15$. Twice the number minus 27 is $2n - 27$. Set equal: $n - 15 = 2n - 27$. Subtract $n$: $-15 = n - 27$. Add $27$: $n = 12$.

\item Perimeter $P = 54$, width $w$, length $\ell = 2w + 3$. Use $P = 2(\ell + w)$: 
\\ $54 = 2((2w+3) + w) = 2(3w + 3) = 6w + 6$. \\ Subtract $6$: $48 = 6w$. So $w = 8$ cm. (Length $= 19$ cm.)

\item "3 more than a number is 20" $\Rightarrow n + 3 = 20$, so $n = 17$. Twice the number $= 2n = 34$.

\item $7x - 3 + 5x + 8 = (7x + 5x) + (-3 + 8) = 12x + 5$.

\item $2(3x - 5) - (x + 7) = 6x - 10 - x - 7 = 5x - 17$.

\item $6 - [2 - (3x - 4)] = 6 - [2 - 3x + 4] = 6 - [6 - 3x] = 6 - 6 + 3x = 3x$.

\item Substitute $y=-2$: $4(-2) - 2(3(-2) - 5) = -8 - 2(-6 - 5) = -8 - 2(-11) = -8 + 22 = 14$.

\item $a=3,\ b=-5$: $2a^2 - ab + b = 2(9) - (3)(-5) + (-5) = 18 + 15 - 5 = 28$.

\item Factor the GCF: $12x + 18 = 6(2x + 3)$.

\item $(3x^2 - 2x + 5) + (x^2 + 7x - 3) = 4x^2 + 5x + 2$.

\item $(5x^2 + 6x - 4) - (2x^2 - x + 9) = (5x^2 - 2x^2) + (6x - (-x)) + (-4 - 9) = 3x^2 + 7x - 13$.

\item $-3x(4x - 7) = -12x^2 + 21x$.

\item $(x + 6)(x - 2) = x^2 - 2x + 6x - 12 = x^2 + 4x - 12$.

\item $(2x - 3)(x + 5) = 2x^2 + 10x - 3x - 15 = 2x^2 + 7x - 15$.

\item $x^2 + 9x + 14$. Find numbers that multiply to $14$ and add to $9$: $2$ and $7$. So $(x + 2)(x + 7)$.

\item $6x^2 + 11x + 3$. Factor by grouping:
\\ Find two numbers with product $6\cdot 3 = 18$ and sum $11$: $9$ and $2$. 
\\ Rewrite middle term: $6x^2 + 9x + 2x + 3$. 
\\ Group: $(6x^2 + 9x) + (2x + 3) = 3x(2x + 3) + 1(2x + 3) = (2x + 3)(3x + 1)$.
\end{enumerate}

\vfill
\hrulefill

\small
\textbf{Teacher notes:} Items 1--7 target one- and two-step equations and simple modeling; 8--13 emphasize structure and evaluation; 14--20 address polynomial operations and factoring.

\end{document}
